\section*{Overview}

A Cartesian tree turns an array into a binary tree that preserves \emph{both} index order and value order:

\begin{itemize}
  \item inorder traversal gives the original array order,
  \item the tree satisfies a heap property on values.
\end{itemize}

That combination makes it useful whenever a range wants to split around its minimum or maximum element.

\section*{When to Use It}

Use it when:

\begin{itemize}
  \item interval structure depends on the minimum or maximum element of the interval,
  \item RMQ should become an LCA problem,
  \item the statement hides recursive splits around dominant elements.
\end{itemize}

\section*{Core Idea}

For a min Cartesian tree, the root of any subarray is its minimum element. The left child is built from the part to the
left, and the right child from the part to the right. A monotonic stack builds the whole tree in linear time.

\section*{Key Insight}

The tree stores the same information as repeated ``pick the minimum and recurse,'' but without paying \(O(n^2)\) for
that recursion literally. The stack simulates how intervals nest around decreasing minima.

\section*{Worked Problem}

\paragraph{Problem.}
Given an array, repeatedly choose the minimum element of a segment as the segment's root and split left/right. Build the
entire recursion tree.

\paragraph{Why Cartesian tree fits.}
That recursion tree \emph{is} the min Cartesian tree. The array index order becomes the inorder traversal, so subtree
boundaries match contiguous ranges automatically.

\section*{Problem Pattern}

This structure appears in:

\begin{itemize}
  \item RMQ to LCA reductions,
  \item histogram problems,
  \item divide-and-conquer recurrences that always pivot on the min or max element.
\end{itemize}

\section*{Correctness Intuition}

The monotonic stack keeps the current right spine of the tree. When a smaller value arrives, larger values to its left
can no longer stay above it in a min-heap tree, so they are popped and attached as its left child or as descendants on
the correct side. The inorder order remains the original index order throughout.

\section*{Complexity Analysis}

Building the tree takes \(O(n)\) time and \(O(n)\) memory.

\section*{Implementation}

The code constructs the min Cartesian tree and returns:

\begin{itemize}
  \item the root index,
  \item parent,
  \item left child,
  \item right child arrays.
\end{itemize}

\section*{Common Pitfalls}

\begin{itemize}
  \item Mixing the min-tree and max-tree inequalities.
  \item Ignoring equal values; the tie rule must be consistent.
  \item Forgetting that the tree shape depends on whether ties favor earlier or later indices.
\end{itemize}

\section*{Variants / Extensions}

\begin{itemize}
  \item Max Cartesian tree.
  \item RMQ via Euler tour + LCA on the Cartesian tree.
  \item Parsing interval DP recurrences through subtree structure.
\end{itemize}

\section*{Practice Problems}

\begin{itemize}
  \item Build the Cartesian tree of an array.
  \item Reduce RMQ to LCA.
  \item Reconstruct recursion around interval minima or maxima.
\end{itemize}

\section*{References}

\begin{itemize}
  \item \href{https://cp-algorithms.com/graph/rmq_linear.html}{cp-algorithms: Solve RMQ by Finding LCA}.
  \item \href{https://en.oi-wiki.org/ds/cartesian-tree/}{OI Wiki: Cartesian Tree}.
  \item \href{https://codeforces.com/catalog}{Codeforces Catalog}.
\end{itemize}
